\scanpage{154}
\mspage{99}
$i_X$ est dans $W_S$, donc iso dans $(\mathcal M_{/S})(W_S)^{-1}=\mathrm{Ho}^{\mathrm{lc}}_W$, et $p$ est donc iso. On retrouve donc $\rho i\simeq\mathrm{id}$, donc que $\rho$ est une \uline{rétraction} de $\mathrm{Ho}^{\mathcal L}_W$ sur $\mathrm{Ho}^{\mathrm{lc}}_W$.

Pour prouver l'adjonction, en procédant comme dans le cas de $\mathrm{Cat}$, on est ramené à prouver que pour tt $X$,
\[
\left\{\begin{aligned}
&C(i_X)=i_{C(X)}:\quad C(X)\rightrightarrows C\deleted{\unsure{C|(}}(C(X))\\
&\text{\uline{dans }}\mathrm{Ho}^{\mathcal L}_W\qquad(\text{loc. par }W_S^{\mathcal L})\text{.}
\end{aligned}\right.
\]

J'y étais arrivé dans $\mathrm{Cat}$, en utilisant un morphisme
\[
C(C(X))\xrightarrow{c}C(X)
\]
(essentiellement, un morphisme de composition des chemins) tel que $c\in W_S^{\mathcal L}$, et que $c\circ C(i_X)=c\circ i_{C(X)}$. Ça avait l'air très délicat, car c'était à grand peine que j'étais arrivé à construire un foncteur $\underline{\mathrm{Ch}}$ (``chemins'') \unsure{,} raisonnable, où on puisse bel et bien composer les chemins ! Je renonce pour le moment à axiomatiser les propriétés du foncteur $\underline{\mathrm{Ch}}$ dont j'ai eu besoin pour faire marcher la machine (je devrais y revenir tantôt).
Je vais commencer plutôt par le ``cas Quillen'', avec des $i_X\in\mathrm{Cof}_0$.
% archive sheet number: 153
